\section{Технический проект}
\subsection{Общая характеристика организации решения задачи}

Необходимо разработать концепцию пользовательского интерфейса для закрытого корпоративного мессенджера, который будет спроектирован как единое, адаптивное, кроссплатформенное веб-приложение, обеспечивающий комфортную работу на различных устройствах (компьютеры, планшеты, мобильные телефоны). В его основе лежит четкая структура и логика, призванная упростить доступ к ключевым функциям.

\begin{enumerate}
	\item Регистрация и авторизация пользователя путем идентификатора (логин) и аутентификатора (пароль) на страницах register и login.
	\item Возможность создавать групповые чаты.
	\item Возможность вести личный чат с другими пользователями. 
	\item Роли пользователей в групповых чатах с различыми правами управления группой.
	\item Поиск сообщений по релевантности и дате.
	\item Просмотр списка участников и их активность (онлайн/оффлайн) в реальном времени.
	\item  Возможность отправлять эмодзи и файлы разных типов.
\end{enumerate}
\subsection{Обоснование выбора технологии проектирования}

Для реализации проектирования интерфейса взаимодействия веб-приложения "Закрытый корпоративный мессенджер" с пользователем были рассмотрены различные подходы и технологии к фронтенд-разработке. Наилучшим стал выбор комбинации языков программирования JavaScript и CSS, а также языка разметки HTML. Этот выбор обусловлен стремлением к достижению максимального контроля над каждым аспектом функциональности, оптимизации производительности и минимизации внешних зависимостей. Такой подход обеспечивает высокую гибкость в процессе разработки, а также эффективность и масштабируемость готового решения. 

\subsection{JavaScript для обеспечения динамического взаимодействия и интерактивности}

Для обеспечения динамического взаимодействия и интерактивности взят нативный JavaScript без использования фреймворков и библиотек. Java-Script файлы (app.js, login.js, register.js) являются "мозгом" клиентской части веб-приложения, обеспечивая всю интерактивность, динамические обновление контента и асинхронное взаимодействие с серверной частью. Одним из важных аспектов выбора данного языка программирования является уникальная способность работать напрямую в браузере. Это позволяет пользователям открывать мессенджер на любом устройстве - смартфон, планшет, ноутбук - без установки дополнительных программ.

\subsubsection {Функциональные возможности, реализованные с помощью JavaScript} 

Одна из самых главных возможностей является - Управление аутентификации и регистрации (login.js, register.js).

Файлы login.js и register.js представляют собой автономные модули, ответственные за логику, связанную с входом и регистрацией пользователей соответственно. Они демонстрируют ключевые принципы нативной JavaScript-разработки: прямое взаимодействие с DOM, асинхронные запросы и управление локальным состоянием без использования фреймворков.

Обработка отправки форм. Скрипты перехватывают событие отправки форм (sumbit) для элементов с id="login-form" и id="register-form". Вместо стандартной отправки страницы, данные из полей ввода "логин" и "пароль" извлекаются и подготавливаются для асинхронной передачи. 

Клиентская валидация данных. До отправки на сервер выполняется проверка их корректности непосредственно в браузере пользователя. Это повышает отзывчивость приложения, так как пользователь получает мгновенную обратную связь без задержек, связанных с сетевым запросом к серверу. Принцип работы:
\begin{enumerate}
	\item Проверка на пустоту. Скрипт проверяет, не являются ли поля username (логин пользователя) и password (пароль пользователя) пустыми.
	\item Если одно из полей пустое, то к соответствующему полю ввода добавляется CSS класс input-error (ошибка ввода). Этот класс изменяет границу поля на красный цвет и фон на светло-красный, визуально сигнализируя пользователю об ошибке.
	\item На пустое поле, функция showToast() вызывается для отображения всплывающего сообщения: "Пожалуйста, заполните это поле".
	\item Выполнение функции прерывается (return), предотвращая отправку некорректных данных на сервер.
	\item Удаление индикации ошибок. При каждом новом событии input (ввода текста в поле) или при попытке повторной отправки формы, ранее добавленный класс input-error удаляется, чтобы очистить визуальную индикацию предыдущих ошибок. Это создает динамичную и понятную обратную связь для пользователя.
\end{enumerate}

Асинхронное взаимодействие с сервером API. После успешной клиентской валидации данные отправляются на сервер без перезагрузки страницы. Это достигается за счет использования современного API fetch(). Принцип работы:

\begin{enumerate}
	\item Создается объект запроса, включающий: URL, метод, заголовки, тело запроса.
	\item URL для входа /login, для регистрации /register.
	\item Метод всегда POST, так как происходит отправка данных.
	\item Заголовки headers, указывается Content-Type: application/json, чтобы сервер знал, что ожидать JSON-данные.
	\item Тело запроса body. Даннные пользователя индикатор (логин, никнейм) и аутентификатор (пароль) конвертируются в JSON-строку с помощью JSON.stringify().
	\item Отправка запроса fetch('/api/login', {...}) или fetch('/api/register', {...}) инициирует сетевой запрос.
	\item Обработка успешного ответа или ошибок. Если response.ok и data.success истинно, это означает успешный вход или регистрацию. Если response.ok ложно или data.success ложно, это указывает на ошибку (например, неверные учетные данные, логин уже занят). В этом случае из data.message извлекается сообщение об ошибке, которое отображается пользователю через showToast().
\end{enumerate}

Для информирования пользователя о результате операций (успех или ошибка) используется функция showToast(). Принцип работы:

\begin{enumerate}
	\item Диначеское создание элементов. Функция showToast(message, type) принимает текст сообщения и его тип success или error.
	\item Она динамически создает новый HTML-элемент div (toastElement), который служит контейнером для уведомления.
	\item К этому элементу добавляются CSS-классы alert, а также alert-success или alert-error в зависимости от типа сообщения. Эти классы, определенные в login.css и register.css, задают внешний вид уведомления: цвет фона, рамки, тени, отступы.
	\item Внутрь toastElement вставляется иконка и текст сообщения.
	\item Размещение в DOM. Созданный toastElement добавляется в специальный контейнер alert-container, который расположен фиксировано в верхней правой части экрана.
	\item Анимация появления и исчезновения.  CSS-класс slideIn, определенный в login.css и register.css череез @keyframes alertSlideIn, обеспечивает плавное появление уведомления сверху.
	\item Через несколько секунд уведомление автоматически удаляется из DOM путем toastElement.remove().
\end{enumerate}

В совокупности, login.js и register.js демонстрируют примеры эффективного и управляемого JavaScript-кода для реализации критически важных фуннкций аутентификации регистрации в веб-приложении "Закрытый корпоративный мессенджер".

Основная логика чата и интерактивность реализованна в app.js. Центральным JavaScript файлом является app.js, который управляет всей динамической частью интерфейса мессенджера после успешного входа пользователя (index.html). Он координирует отображение данных, реакцию на действия пользователя и взаимодействие с серверным API. 
При загрузке index.html или при перенаправлении на него после входа, app.js выполняет начальную настройку:

\begin{enumerate}
	\item Проверка аутентификации. Прогр

\subsection{Диаграмма компонентов базы данных}

Диаграмма компонентов базы данных представляет собой визуализацию структуры и взаимосвязей основных модулей системы хранения данных.

\begin{figure}[ht]
\center{\includegraphics[width=1\linewidth]{ЕР модель}}
\caption{Диаграмма компонентов базы данных}
\label{comp:image}
\end{figure}

ER-модель отражает строгую структуру данных, построенную вокруг взаимодействия пользователей в групповых и приватных чатах. В центре системы находится сущность "Пользователь", хранящая учетные данные и выступающая отправителем сообщений. Группы организуют коллективное общение, фиксируя создателя и время создания, а Сообщения делятся на групповые и приватные, с обязательной привязкой к отправителю и временной меткой.

Система ролей реализована через связующую сущность "GroupMember", которая определяет права участников групп (владелец, администратор, участник) и дату их присоединения. Это обеспечивает гибкое управление доступом внутри чатов. Для работы с файлами используется сущность "Attachment", привязанная к сообщениям и содержащая технические данные о файлах (путь, тип, оригинальное имя), что позволяет единообразно обрабатывать вложения как в групповых, так и в личных сообщениях.

Модель подчеркивает важность целостности данных через систему внешних ключей, гарантируя, что сообщения не могут существовать без отправителя или чата. Поле "timestamp" обеспечивает учет активности пользователей, а минималистичная реализация приватных чатов (прямая связь отправитель-получатель) упрощает архитектуру. Вынос вложений в отдельную таблицу позволяет менять способ хранения файлов без изменения структуры сообщений.

\subsection{Таблицы базы данных}
\subsubsection{User - Пользователи системы}
\begin{xltabular}{\textwidth}{|l|l|X|}
	\caption{Атрибуты сущности "Пользователи"\label{user:table}}\\ \hline
	\centrow Поле & \centrow Тип & \centrow Описание \\ \hline
	\thead{id} & \thead{INTEGER} & Первичный ключ (PK), уникальный идентификатор пользователя \\ \hline
	\thead{username} & \thead{TEXT} & Имя пользователя (уникальное) \\ \hline
	\thead{password} & \thead{TEXT} & Хэшированный пароль пользователя \\ \hline
\end{xltabular}

Таблица "User" хранит информацию о зарегистрированных пользователях мессенджера. Каждый пользователь имеет уникальный идентификатор, имя пользователя и пароль. Является центральной сущностью системы, с которой связаны все остальные таблицы.

\subsubsection{Group - Групповые чаты}
\begin{xltabular}{\textwidth}{|l|l|X|}
	\caption{Атрибуты сущности "Групповые чаты"\label{group:table}}\\ \hline
	\centrow Поле & \centrow Тип & \centrow Описание \\ \hline
	\thead{group\_id} & \thead{INTEGER} & Первичный ключ (PK), ID группы \\ \hline
	\thead{name} & \thead{TEXT} & Название группы \\ \hline
	\thead{creator\_id} & \thead{INTEGER} & Внешний ключ (FK→User), ID создателя группы \\ \hline
	\thead{created\_at} & \thead{INTEGER} & Timestamp создания группы \\ \hline
\end{xltabular}

Таблица "Group" содержит информацию о групповых чатах. Каждая группа имеет уникальный идентификатор, название, создателя и время создания. Создатель группы автоматически становится владельцем (owner).

\subsubsection{GroupMember - Участники групп с ролями}
\begin{xltabular}{\textwidth}{|l|l|X|}
	\caption{Атрибуты сущности "Участники групп"\label{groupmember:table}}\\ \hline
	\centrow Поле & \centrow Тип & \centrow Описание \\ \hline
	\thead{group\_id} & \thead{INTEGER} & Внешний ключ (FK→Group), ID группы \\ \hline
	\thead{user\_id} & \thead{INTEGER} & Внешний ключ (FK→User), ID пользователя \\ \hline
	\thead{role} & \thead{TEXT} & Роль в группе: 'owner' (владелец), 'admin' (администратор), 'member' (участник) \\ \hline
	\thead{joined\_at} & \thead{INTEGER} & Timestamp вступления пользователя в группу \\ \hline
\end{xltabular}

Таблица "GroupMember" определяет состав участников групп и их роли. Реализует систему прав доступа, аналогичную Telegram, где владелец имеет максимальные права, администраторы - ограниченные права управления, а участники - базовые права.

\subsubsection{GroupMessage - Сообщения в группах}
\begin{xltabular}{\textwidth}{|l|l|X|}
	\caption{Атрибуты сущности "Групповые сообщения"\label{groupmessage:table}}\\ \hline
	\centrow Поле & \centrow Тип & \centrow Описание \\ \hline
	\thead{message\_id} & \thead{INTEGER} & Первичный ключ (PK), ID сообщения \\ \hline
	\thead{group\_id} & \thead{INTEGER} & Внешний ключ (FK→Group), ID группы \\ \hline
	\thead{user\_id} & \thead{INTEGER} & Внешний ключ (FK→User), ID отправителя \\ \hline
	\thead{message\_text} & \thead{TEXT} & Текст сообщения (может быть пустым для сообщений с вложениями) \\ \hline
	\thead{timestamp} & \thead{INTEGER} & Время отправки сообщения (timestamp) \\ \hline
\end{xltabular}

Таблица "GroupMessage" хранит все сообщения, отправленные в групповых чатах. Каждое сообщение связано с группой и пользователем-отправителем. Поддерживает текстовые сообщения и сообщения с вложениями (через таблицу Attachment).

\subsubsection{PrivateMessage - Личные сообщения}
\begin{xltabular}{\textwidth}{|l|l|X|}
	\caption{Атрибуты сущности "Личные сообщения"\label{privatemessage:table}}\\ \hline
	\centrow Поле & \centrow Тип & \centrow Описание \\hline
	\thead{id} & \thead{INTEGER} & Первичный ключ (PK), уникальный идентификатор сообщения \\ \hline
	\thead{sender\_id} & \thead{INTEGER} & Внешний ключ (FK→User), ID отправителя сообщения \\ \hline
	\thead{receiver\_id} & \thead{INTEGER} & Внешний ключ (FK→User), ID получателя сообщения \\ \hline
	\thead{message\_text} & \thead{TEXT} & Текст сообщения (может быть пустым для сообщений с вложениями) \\ \hline
	\thead{timestamp} & \thead{INTEGER} & Время отправки сообщения (timestamp) \\ \hline
\end{xltabular}

Таблица "PrivateMessage" хранит историю приватных сообщений между пользователями. Каждое сообщение связывает отправителя и получателя, поддерживает текстовый контент и вложения (через таблицу Attachment). Реализует функционал личной переписки между сотрудниками организации.

\subsubsection{Attachment - Вложения к сообщениям}
\begin{xltabular}{\textwidth}{|l|l|X|}
	\caption{Атрибуты сущности "Вложения"\label{attachment:table}}\\ \hline
	\centrow Поле & \centrow Тип & \centrow Описание \\ \hline
	\thead{id} & \thead{INTEGER} & Первичный ключ (PK), уникальный идентификатор вложения \\ \hline
	\thead{message\_type} & \thead{TEXT} & Тип сообщения: 'general' (системные), 'group' (групповые), 'private' (личные) \\ \hline
	\thead{message\_id} & \thead{INTEGER} & ID сообщения, к которому относится вложение \\ \hline
	\thead{file\_path} & \thead{TEXT} & Путь к файлу в файловом хранилище системы \\ \hline
	\thead{mime\_type} & \thead{TEXT} & MIME-тип файла для корректной обработки клиентом \\ \hline
	\thead{filename} & \thead{TEXT} & Оригинальное имя файла для отображения пользователю \\ \hline
\end{xltabular}

Таблица "Attachment" обеспечивает хранение файловых вложений для всех типов сообщений в системе. Универсальная структура позволяет прикреплять файлы к групповым чатам, личным сообщениям и системным уведомлениям. Содержит метаданные файлов для правильного отображения в клиентских приложениях.

\subsection{Пакет models}

Перенос всех запросов к базе данных из views в models обеспечивает достиженение более чистой и масштабируемой архитектуры приложения. Так, views отвечают исключительно за обработку HTTP-запросов, валидацию входных данных и формирование ответов, в то время как вся работа с базой данных сосредоточена в models.

Подход обеспечивает чёткое разделение ответственности между компонентами системы. Модели инкапсулируют бизнес-логику, связанную с данными, включая сложные проверки прав доступа, транзакции и обеспечение целостности данных. Это делает код более безопасным, так как все SQL-запросы выполняются в одном контролируемом слое, где можно централизованно применять меры защиты от инъекций и других уязвимостей.

\subsubsection{UserModel.py - управление пользователями}
Модуль обеспечивает взаимодействие с таблицей User базы данных. Основные функции:
\begin{itemize}
	\item аутентификация через метод authenticate(), который проверяет соответствие хэша пароля в базе данных;
	\item создание новых пользователей с валидацией данных и хэшированием паролей;
	\item поиск пользователей по имени или ID для системы контактов.
\end{itemize}

Запросы используют параметризованные SQL-выражения для защиты от инъекций. Хэши паролей хранятся с использованием алгоритма bcrypt.

\subsubsection{GroupModel.py - управление группами}
Модуль работает с таблицами Group и GroupMember, реализуя сложную логику управления чатами:
\begin{itemize}
	\item создание групп с автоматическим назначением создателя владельцем;
	\item система ролей (владелец/администратор/участник) с проверкой прав доступа;
	\item механизмы добавления/исключения участников с валидацией полномочий;
	\item управление составом групп через транзакционные операции.
\end{itemize}

Особенность реализации — каскадные проверки прав при любых изменениях состава групп.

\subsubsection{MessageModel.py - работа с сообщениями}
Модуль обслуживает таблицы GroupMessage, PrivateMessage и Attachment, предоставляя:
\begin{itemize}
	\item единый интерфейс для работы с разными типами сообщений;
	\item поддержку вложений с обработкой MIME-типов;
	\item поиск по истории сообщений с пагинацией;
	\item механизмы редактирования и удаления с проверкой прав.
\end{itemize}

Операции с сообщениями реализованы через транзакции для сохранения целостности данных. При удалении сообщений автоматически удаляются связанные вложения.

\subsubsection{Принципы работы с базой данных}
\begin{itemize}
	\item запросы используют контекстные менеджеры для управления соединениями;
	\item ошибки логируются и обрабатываются на уровне моделей;
	\item данные возвращаются в виде словарей/объектов, а не сырых SQL-результатов;
	\item реализовано единое подключение через utils.get\_db\_cursor().
\end{itemize}

Вынесенные к базе данных запросы в отдельный пакет облегчают возможный переход на другую СУБД или использование ORM. Поскольку все запросы сосредоточены в одном месте, их модификация при смене технологии работы с данными потребует минимальных изменений в остальных частях приложения.

Всё это делает систему более гибкой, легкочитаемой и готовой к возможным изменениям.

\subsection{Подключение к базе данных}

Механизм взаимодействия приложения с базой данных реализован с помощью контекстных менеджеров, которые обеспечивают безопасное управление ресурсами. В модуле db\_utils.py определены два ключевых метода: get\_db\_connection() и get\_db\_cursor(). Первый предоставляет соединение с базой данных, второй — курсор для выполнения SQL-запросов.

Использование контекстных менеджеров обеспечивает автоматическое закрытие соединения и курсора после завершения операций с таблицами, что минимизирует вероятность утечек ресурсов. Такой подход делает код более лаконичным и безопасным, а также исключает риск возникновения состояния «зависших» соединений, что особенно важно при высокой нагрузке на систему.

\subsection{Хеширование паролей}

Для хранения паролей пользователей реализована схема хеширования с использованием библиотеки bcrypt. В модуле pswd\_utils.py метод hash\_password() выполняет хеширование, а check\_password() — проверку соответствия введенного пароля его хешу.

Алгоритм bcrypt предоставляет высокий уровень безопасности за счет использования адаптивного хеширования, усложняющего процесс взлома методом перебора. Благодаря этому защита учетных данных пользователей соответствует современным требованиям информационной безопасности, а аутентификация в системе надежно защищена от нежелательных атак.

\subsection{Диаграмма последовательности}

На рисунке \ref{data:image} представлена схема обмена данными между сценариями компонента при вызове компонента на странице сайта.

\begin{figure}[H]
	\center{\includegraphics[width=1\linewidth]{Диаграмма последовательности}}
	\caption{Диаграмма компонентов}
	\label{data:image}
\end{figure}

\begin{enumerate}[leftmargin=*,label=\textbf{\arabic*.}]
	\item \textbf{Загрузка страницы (инициализация)}
	\begin{itemize}
		\item пользователь открывает в браузере главную страницу (index.html);
		\item браузер отправляет HTTP GET-запрос на сервер;
		\item серверный скрипт run.py (Waitress) получает запрос;
		\item запрос передаётся в WSGI-приложение (app.py);
		\item маршрутизатор (routes.py) определяет обработчик IndexView;
		\item пользователю возвращается HTML-страница и статические файлы (CSS/JS).
	\end{itemize}
	
	\item \textbf{Процесс аутентификации}
	\begin{itemize}
		\item пользователь вводит логин и пароль в форму входа;
		\item браузер отправляет POST /login запрос (JSON);
		\item LoginView проверяет учётные данные:
		\begin{itemize}
			\item через utils.py выполняется SQL-запрос к таблице users;
			\item пароль проверяется с помощью bcrypt.
		\end{itemize}
		\item при успешной аутентификации:
		\begin{itemize}
			\item устанавливается cookie с user\_id;
			\item происходит редирект на главную страницу.
		\end{itemize}
	\end{itemize}
	
	\item \textbf{Отправка нового сообщения}
	\begin{itemize}
		\item пользователь вводит текст и/или прикрепляет файлы;
		\item браузер формирует FormData и отправляет POST /send\_message;
		\item SendMessageView выполняет:
		\begin{itemize}
			\item проверку прав через cookies;
			\item сохранение файлов в static/uploads;
			\item запись сообщения в БД (таблицы messages или group\_messages);
		\end{itemize}
		\item клиент получает JSON с результатом операции;
	\end{itemize}
	
	\item \textbf{Получение сообщений (реализация чата)};
	\begin{itemize}
		\item браузер периодически опрашивает сервер (GET /get\_messages);
		\item GetMessageView запрашивает новые сообщения из БД;
		\item сервер возвращает JSON с массивом сообщений;
		\item браузер динамически обновляет интерфейс чата;
	\end{itemize}
	
	\item \textbf{Особенности работы}
	\begin{itemize}
		\item все запросы проходят через единую точку входа (app.py);
		\item маршрутизатор выбирает соответствующий View-класс;
		\item бизнес-логика сосредоточена во views;
		\item работа с БД вынесена в utils.py;
		\item файлы хранятся локально в файловой системе;
		\item состояние сессии поддерживается через cookies.
	\end{itemize}
\end{enumerate}

\subsection{Диаграмма размещения}

Диаграмма размещения (рис.~\ref{place:image}) отражает физические взаимосвязи между программными и аппаратными компонентами системы.

\begin{figure}[ht]
\center{\includegraphics[width=1\linewidth]{Диаграмма развёртывания}}
\caption{Диаграмма развёртывания}
\label{place:image}
\end{figure}

Диаграмма развертывания отображает архитектуру корпоративного мессенджера. На стороне клиента взаимодействие с системой осуществляется через веб-браузер, который загружает статические файлы интерфейса — HTML-шаблоны, JavaScript-код и CSS-стили.

Серверная часть реализована на Python с использованием WSGI-сервера Waitress, обрабатывающего входящие HTTP-запросы на порту 8000. Логика приложения организована в виде пакетов views и models. Пакет views содержит специализированные модули для обработки запросов: auth.py отвечает за аутентификацию, groups.py управляет групповыми чатами, message.py обрабатывает сообщения, p\_chat.py обеспечивает личную переписку. Пакет models включает UserModel.py, GroupModel.py и MessageModel.py, инкапсулирующие работу с данными. Вспомогательные функции вынесены в utils.py и mimes.py.

Хранение данных организовано в двух форматах. Реляционная база SQLite (файл data.db) содержит таблицы User, Group, GroupMember, GroupMessage, PrivateMessage и Attachment, обеспечивающие хранение структурированной информации о пользователях, чатах, сообщениях и вложениях. Файловое хранилище в папках static/uploads и templates обслуживает медиафайлы и шаблоны интерфейса.

Взаимодействие между компонентами осуществляется через API между клиентом и сервером. WSGI-сервер передает запросы в Python-приложение, где они обрабатываются соответствующими модулями пакета views. Для работы с данными views используют методы пакета models, которые выполняют SQL-запросы к базе данных. Файловые вложения передаются напрямую через специальные обработчики.

Архитектура системы сохраняет принцип самодостаточности — все компоненты развертываются на едином сервере без внешних зависимостей. Четкое разделение на модули views и models обеспечивает безопасность и упрощает дальнейшее развитие системы, сохраняя при этом возможность автономной работы в корпоративной сети.

